% !TEX program = xelatex
\documentclass[11pt,a4paper]{article}
\usepackage{algo-preamble}

\title{\textbf{L14 \textperiodcentered{} Δυναμικός Προγραμματισμός I}\\[2pt]
\large Memoization \& Σταθμισμένος Χρονοπρογραμματισμός}
\author{Algorithms Class Hub}
\date{\small Αλγόριθμοι και Πολυπλοκότητα (K17) \textperiodcentered{} ΕΚΠΑ}

\begin{document}
\maketitle

\begin{center}\itshape\small
Η ιδέα του δυναμικού προγραμματισμού: επικαλυπτόμενα υποπροβλήματα και
memoization. Fibonacci, σταθμισμένος χρονοπρογραμματισμός διαστημάτων, και η
συνταγή 4 βημάτων του DP.
\end{center}

\section{Τι θα δούμε}

Ξεκινάει το τελευταίο μεγάλο κεφάλαιο: ο \textbf{δυναμικός προγραμματισμός}
(dynamic programming, DP). Είναι το πιο εξεταζόμενο κομμάτι της ύλης --- και αυτό
που οι φοιτητές βρίσκουν πιο «μαγικό». Στόχος αυτής της διάλεξης: να πάψει να
είναι μαγικό.

Θα δούμε την ιδέα μέσα από το πιο απλό παράδειγμα (Fibonacci), μετά το πρώτο
«πραγματικό» πρόβλημα (\textbf{σταθμισμένος χρονοπρογραμματισμός διαστημάτων}),
και θα κωδικοποιήσουμε μια \textbf{συνταγή 4 βημάτων} που εφαρμόζεται σε
\textbf{κάθε} DP πρόβλημα.

\section{Τα τρία αλγοριθμικά μοντέλα}

\begin{keypoint}
\begin{itemize}
  \item \textbf{Απληστία} --- χτίσε τη λύση σταδιακά, βελτιστοποιώντας μυωπικά
  ένα τοπικό κριτήριο.
  \item \textbf{Διαίρει και κυρίευε} --- διάσπασε σε \textbf{ανεξάρτητα}
  υποπροβλήματα, λύσε το καθένα ξεχωριστά, συνδύασε.
  \item \textbf{Δυναμικός προγραμματισμός} --- διάσπασε σε μια σειρά από
  \textbf{επικαλυπτόμενα} υποπροβλήματα, και δόμησε σωστές λύσεις για όλο και
  μεγαλύτερα υποπροβλήματα.
\end{itemize}
\end{keypoint}

\begin{intuition}
\textbf{Η κρίσιμη λέξη είναι «επικαλυπτόμενα».} Στο διαίρει-και-κυρίευε τα δύο
μισά ενός πίνακα είναι \textbf{ξένα} --- δεν μοιράζονται τίποτα. Στον δυναμικό
προγραμματισμό τα υποπροβλήματα \textbf{ξανασυναντιούνται}: το ίδιο υποπρόβλημα
χρειάζεται πολλές φορές. Αν το λύνεις ξανά κάθε φορά, πληρώνεις εκθετικά. Αν το
λύνεις \textbf{μία φορά και το θυμάσαι}, πληρώνεις γραμμικά. Αυτή είναι όλη η
ιδέα.
\end{intuition}

\begin{notebox}
\textbf{Από πού το όνομα;} Ο Richard Bellman θεμελίωσε τον δυναμικό
προγραμματισμό τη δεκαετία του 1950. «Programming» εδώ σημαίνει
\textbf{σχεδιασμός} (όπως στο «linear programming»), όχι κώδικας. Ο Bellman
διάλεξε επίτηδες ένα εντυπωσιακό όνομα --- κάτι που ούτε ένα μέλος του
Κογκρέσου δεν θα μπορούσε να απορρίψει ως χρηματοδότηση. Διάσημοι DP αλγόριθμοι:
\textbf{Viterbi}, \textbf{unix diff}, \textbf{Smith-Waterman},
\textbf{Bellman-Ford}, \textbf{CKY}.
\end{notebox}

\section{Το παράδειγμα-εκκίνηση: Fibonacci}

Η ακολουθία Fibonacci ορίζεται αναδρομικά:
\[
F(n) = F(n-1) + F(n-2), \qquad F(0) = 0,\ F(1) = 1.
\]
Η απευθείας μεταφορά σε κώδικα:

\begin{codebox}
\begin{verbatim}
fib(n):
  Εάν n <= 1: επίστρεψε n
  επίστρεψε fib(n-1) + fib(n-2)
\end{verbatim}
\end{codebox}

Κομψό --- και \textbf{καταστροφικό}. Η πολυπλοκότητα είναι $\approx O(2^n)$,
εκθετική.

\subsection{Γιατί είναι τόσο αργό; Επικάλυψη υποπροβλημάτων}

Όταν υπολογίζουμε το $F(5)$ με την αφελή αναδρομή, το δέντρο κλήσεων υπολογίζει
το $F(3)$ \textbf{2 φορές} και το $F(2)$ \textbf{3 φορές}. Όσο μεγαλώνει το $n$,
το δέντρο αναδρομής \textbf{εκρήγνυται} --- επειδή τα ίδια υποπροβλήματα
υπολογίζονται ξανά και ξανά.

\subsection{Η λύση: memoization}

\begin{keypoint}
\textbf{Απομνημόνευση (memoization):} κάθε φορά που υπολογίζεις έναν όρο $F(k)$,
\textbf{αποθήκευσέ τον} σε έναν πίνακα. Όταν τον ξαναζητήσει η αναδρομή,
\textbf{επίστρεψε την αποθηκευμένη τιμή} αντί να ξαναξεκινήσεις τον υπολογισμό.
\end{keypoint}

\begin{codebox}
\begin{verbatim}
fib2(n):
  Εάν n <= 1: επίστρεψε n
  A[0] <- 0 ;  A[1] <- 1
  Για i = 2 έως n:
    A[i] <- A[i-1] + A[i-2]
  επίστρεψε A[n]
\end{verbatim}
\end{codebox}

Κάθε $F(k)$ υπολογίζεται \textbf{ακριβώς μία φορά}. Η πολυπλοκότητα πέφτει από
εκθετική $O(2^n)$ σε \textbf{γραμμική $O(n)$}.

\section{Σταθμισμένος Χρονοπρογραμματισμός Διαστημάτων}

Τώρα το πρώτο πραγματικό πρόβλημα. Είναι το πρόβλημα του χρονοπρογραμματισμού
διαστημάτων από το μάθημα των άπληστων αλγορίθμων --- \textbf{με μια προσθήκη}.

\textbf{Είσοδος:} $n$ αιτήματα. Το αίτημα $j$ ξεκινά τη στιγμή $s_j$, τελειώνει
τη στιγμή $f_j$, και έχει \textbf{βαρύτητα} $v_j$. Δύο αιτήματα είναι
\textbf{συμβατά} αν δεν επικαλύπτονται.

\textbf{Στόχος:} υποσύνολο $S$ ανά δύο συμβατών αιτημάτων με \textbf{μέγιστο
άθροισμα βαρυτήτων} $\sum_{i \in S} v_i$.

\begin{warningbox}
\textbf{Ο άπληστος αλγόριθμος αποτυγχάνει θεαματικά.} Στους άπληστους
αλγορίθμους, με όλες τις βαρύτητες ίσες με $1$, το «μικρότερος χρόνος λήξης»
έδινε βέλτιστη λύση. Με αυθαίρετες βαρύτητες όμως, μια εργασία που τελειώνει
νωρίς αλλά αξίζει $1$ μπορεί να μπλοκάρει μία που αξίζει $100$. Κανένα απλό
κριτήριο σειράς δεν αρκεί --- χρειαζόμαστε DP.
\end{warningbox}

\subsection{Η προετοιμασία: ταξινόμηση και $p(j)$}

Διατάσσουμε τα αιτήματα κατά \textbf{χρόνο λήξης}:
$f_1 \le f_2 \le \dots \le f_n$. Ορίζουμε:

\begin{keypoint}
$p(j) = $ ο \textbf{μεγαλύτερος} δείκτης $i < j$ τέτοιος ώστε το αίτημα $i$ να
είναι \textbf{συμβατό} με το $j$ (δηλαδή $f_i \le s_j$). Αν κανένα προηγούμενο
αίτημα δεν είναι συμβατό, $p(j) = 0$.
\end{keypoint}

Διαισθητικά: «αν διαλέξω το $j$, ποιο είναι το \textbf{τελευταίο} προηγούμενο
αίτημα που μπορώ να συνδυάσω μαζί του;»

\subsection{Η αναδρομή --- δυαδική επιλογή}

Ορίζουμε $\text{OPT}(j) = $ η τιμή της βέλτιστης λύσης για τα αιτήματα
$\{1, 2, \dots, j\}$. Σκεφτόμαστε το αίτημα $j$: \textbf{μέσα ή έξω;}
\begin{itemize}
  \item \textbf{Το $j$ είναι ΜΕΣΑ στη βέλτιστη λύση.} Τότε αποκλείονται όλα τα
  ασύμβατα με το $j$ --- δηλαδή τα $\{p(j)+1, \dots, j-1\}$. Ό,τι μένει είναι η
  βέλτιστη λύση για τα $\{1, \dots, p(j)\}$. Συνεισφορά:
  $v_j + \text{OPT}(p(j))$.
  \item \textbf{Το $j$ είναι ΕΞΩ.} Τότε η βέλτιστη λύση είναι ακριβώς η βέλτιστη
  για τα $\{1, \dots, j-1\}$. Συνεισφορά: $\text{OPT}(j-1)$.
\end{itemize}
Δεν ξέρουμε ποια περίπτωση ισχύει --- οπότε παίρνουμε τη \textbf{μεγαλύτερη}:

\begin{keypoint}
\[
\text{OPT}(j) =
\begin{cases}
0 & j = 0 \\[4pt]
\max\bigl\{\, v_j + \text{OPT}(p(j)),\ \ \text{OPT}(j-1) \,\bigr\} & j \ge 1
\end{cases}
\]
\end{keypoint}

\subsection{Από την ωμή βία στο memoization}

Η απευθείας αναδρομή \texttt{ComputeOPT(j)} είναι \textbf{εκθετική} --- ακριβώς
όπως ο αφελής Fibonacci, το πλήθος αναδρομικών κλήσεων μεγαλώνει σαν
$T(n) = T(n-1) + T(n-2)$. Η θεραπεία είναι η ίδια: \textbf{απομνημόνευση}.

\begin{codebox}
\begin{verbatim}
Αρχικοποίησε M[1..n] = κενό ;  M[0] = 0

M-ComputeOPT(j):
  Εάν M[j] = κενό:
    M[j] <- max{ v_j + M-ComputeOPT(p(j)),  M-ComputeOPT(j-1) }
  επίστρεψε M[j]
\end{verbatim}
\end{codebox}

\begin{keypoint}
\textbf{Χρόνος εκτέλεσης: $O(n \log n)$.}
\begin{itemize}
  \item Ταξινόμηση κατά χρόνο λήξης: $O(n \log n)$.
  \item Υπολογισμός όλων των $p(j)$: $O(n)$ μετά από ταξινόμηση και κατά χρόνο
  έναρξης.
  \item \texttt{M-ComputeOPT}: κάθε κλήση είναι $O(1)$· κάθε νέα συμπλήρωση μιας
  θέσης $M[j]$ κάνει το πολύ $2$ αναδρομικές κλήσεις. Μέτρο προόδου $\Phi = $
  πλήθος μη-κενών θέσεων του $M[\cdot]$· ξεκινά $0$, φτάνει $\le n$, και κάθε
  «νέα» κλήση το αυξάνει κατά $1$. Άρα συνολικά $O(n)$ κλήσεις.
\end{itemize}
\end{keypoint}

Ισοδύναμα, \textbf{bottom-up} --- γεμίζουμε τον πίνακα με αύξουσα σειρά, χωρίς
αναδρομή:

\begin{codebox}
\begin{verbatim}
M[0] <- 0
Για j = 1 έως n:
  M[j] <- max{ v_j + M[p(j)],  M[j-1] }
\end{verbatim}
\end{codebox}

\subsection{Εύρεση της λύσης}

Ο αλγόριθμος υπολογίζει τη \textbf{μέγιστη τιμή}. Για να βρούμε \textbf{ποια
αιτήματα} την πετυχαίνουν, κάνουμε ένα πέρασμα προς τα πίσω, ρωτώντας σε κάθε $j$
«ποια από τις δύο περιπτώσεις κέρδισε;»

\begin{codebox}
\begin{verbatim}
Βρες-Λύση(j):
  Εάν j = 0: μην εκτυπώσεις τίποτα
  Αλλιώς εάν v_j + M[p(j)] > M[j-1]:    // το j ήταν ΜΕΣΑ
    εκτύπωσε j
    Βρες-Λύση(p(j))
  Αλλιώς:                               // το j ήταν ΕΞΩ
    Βρες-Λύση(j-1)
\end{verbatim}
\end{codebox}

Το πολύ $n$ κλήσεις, άρα $O(n)$.

\section{Η συνταγή του δυναμικού προγραμματισμού}

Κάθε πρόβλημα DP --- όσα θα δούμε στα επόμενα μαθήματα --- λύνεται με τα
\textbf{ίδια 4 βήματα}:

\begin{keypoint}
\begin{enumerate}
  \item \textbf{Χαρακτηρισμός της δομής} του προβλήματος --- όρισε τι είναι ένα
  «υποπρόβλημα» (εδώ: «τα πρώτα $j$ αιτήματα»).
  \item \textbf{Αναδρομικός ορισμός} της τιμής της βέλτιστης λύσης --- η εξίσωση
  $\text{OPT}(\cdot)$.
  \item \textbf{Υπολογισμός} της τιμής της βέλτιστης λύσης --- γέμισε τον πίνακα
  (memoization ή bottom-up).
  \item \textbf{Κατασκευή} της βέλτιστης λύσης από την αποθηκευμένη πληροφορία
  --- το πέρασμα προς τα πίσω.
\end{enumerate}
\end{keypoint}

\begin{intuition}
\textbf{Το πιο δύσκολο βήμα είναι το 1ο.} Μόλις βρεις τη σωστή
«παραμετροποίηση» των υποπροβλημάτων, η αναδρομή σχεδόν γράφεται μόνη της. Σε
επόμενο μάθημα θα δούμε πρόβλημα (knapsack) όπου η πρώτη, «προφανής»
παραμετροποίηση \textbf{δεν δουλεύει} --- και χρειάζεται μια δεύτερη μεταβλητή.
\end{intuition}

\begin{recapbox}
\textbf{Τι κρατάμε από αυτή τη διάλεξη}
\begin{enumerate}
  \item \textbf{Δυναμικός προγραμματισμός} --- διάσπαση σε \textbf{επικαλυπτόμενα}
  υποπροβλήματα· λύσε καθένα μία φορά, αποθήκευσέ το.
  \item \textbf{Fibonacci} --- αφελής αναδρομή $O(2^n)$ λόγω επικάλυψης·
  \textbf{memoization} $\to O(n)$.
  \item \textbf{Σταθμισμένος χρονοπρογραμματισμός} --- μέγιστη βαρύτητα συμβατών
  αιτημάτων. Ο άπληστος αποτυγχάνει.
  \item $p(j) = $ τελευταίο συμβατό προηγούμενο αίτημα. Αναδρομή:
  $\text{OPT}(j) = \max\{v_j + \text{OPT}(p(j)),\ \text{OPT}(j-1)\}$.
  \item Με memoization / bottom-up $\to O(n \log n)$· εύρεση λύσης με πέρασμα
  προς τα πίσω $\to O(n)$.
  \item \textbf{Η συνταγή DP (4 βήματα):} δομή $\to$ αναδρομικός ορισμός $\to$
  υπολογισμός τιμής $\to$ κατασκευή λύσης.
\end{enumerate}
\end{recapbox}

\lecturefooter
\end{document}
